
When a large language model is fine-tuned for a downstream classification task using low-rank adaptation (LoRA), its attention patterns reorganize around task-discriminative tokens. KV cache eviction policies, such as Locret \citep{Huangetal2024}, are trained independently on language modeling loss and learn to prioritize next-token-predictive tokens. When both components are deployed together under a fixed KV budget, the eviction policy discards tokens according to a criterion that was not informed by the task adaptation objective.

This work provides the first direct empirical measurement of this misalignment for a task-fine-tuned LoRA model. The Spearman rank correlation between LoRA-modified attention scores and Locret contextual importance scores is 0.3662 ± 0.0759 across 100 MNLI validation examples on LLaMA-3.1-8B. Every example falls below the 0.70 alignment threshold. This result is consistent with the intuition that task classification and next-token prediction reward different token patterns, but the magnitude of misalignment (ρ ≈ 0.37, indicating approximately 14\% shared variance) had not been quantified before.

The standard pipeline treats fine-tuning and KV compression as independent modules: fine-tune the model with PEFT, then apply a separately-trained compression policy. No prior work has jointly optimized the adaptation and eviction objectives for task-specific performance. The closest related work, arXiv:2604.21335 [Jiang \& Wang, 2026], combines LoRA routing and KV value-group routing but uses language modeling loss throughout, leaving task-specific alignment unaddressed.

JointLoRA-KV is a training procedure designed to close this alignment gap. LoRA adapter parameters and Locret retaining head parameters are optimized together in a single backward pass via the task classification loss. The eviction boundary is made differentiable during training through a sigmoid relaxation of the KV budget mask; an STE bridges the resulting training-inference discrepancy.

The paper reports results from three executed sub-experiments (H-E1, H-M1, H-M2) out of five planned (H-M3, H-M4 not executed). The main findings are: (1) misalignment is confirmed at substantial magnitude; (2) the joint training mechanism is functional at PoC scale, with measurable but sub-threshold improvement over frozen-Locret baseline; (3) joint training is numerically stable across seeds on a small model. The primary claim regarding improvement over sequential fine-tuning (B3) on long-context tasks remains untested.

\textbf{Contributions:}

\begin{enumerate}
\item The first empirical quantification of misalignment between task-adapted LoRA attention patterns and LM-trained Locret eviction scores on a downstream NLI task: Spearman ρ = 0.3662 ± 0.076 across 100 MNLI examples on LLaMA-3.1-8B.
\end{enumerate}

\begin{enumerate}
\item JointLoRA-KV: a joint training procedure for LoRA adapters and KV eviction heads using a differentiable soft KV budget mask with STE, confirmed functional at PoC scale (+1.50pp GLUE over frozen-Locret baseline, formal gate: PARTIAL).
\end{enumerate}

\begin{enumerate}
\item Empirical confirmation that joint training of LoRA and Locret parameters is numerically stable across 3 seeds on a PoC model (d=64, 2 layers; 0 NaN/divergence events). Full-scale stability on LLaMA-3.1-8B was not empirically verified.
\end{enumerate}

